\documentclass[aspectratio=169]{beamer}
\setbeamerfont{headline}{size=\footnotesize}

\mode<presentation>{
    \usetheme{Montpellier}
    \setbeamercovered{transparent}
    \usecolortheme{beaver}
}

\usepackage[english,brazil]{babel}
\usepackage[utf8]{inputenc}
\usepackage{fontawesome}
\usepackage{graphicx}
\usepackage{url}
\usepackage{colortbl}
\usepackage{caption}
\usepackage{booktabs}
\usepackage{multirow}
\usepackage{array}
\usepackage{scalefnt}

\definecolor{mineblue}{RGB}{0,90,180}
\definecolor{mineorange}{RGB}{230,115,0}
\definecolor{minegreen}{RGB}{34,139,34}
\definecolor{minered}{RGB}{180,20,40}
\definecolor{minegray}{RGB}{100,100,100}

\usepackage{tikz}
\usetikzlibrary{arrows.meta,shapes.geometric,positioning,fit,backgrounds,calc,shadows.blur,trees}
\usepackage{pgfplots}
\pgfplotsset{compat=1.18}

\newcommand{\reflexao}[1]{\begin{alertblock}{Reflexão}#1\end{alertblock}}
\newcommand{\key}[1]{\underline{#1}}

\title[Mineração de Dados]{Mineração de Dados: Conceitos e Aplicações}
\subtitle{Módulo 3: Técnicas de Classificação}
\author[Elton Sarmanho]{Prof. Elton Sarmanho\inst{1}\\\footnotesize \href{mailto:eltonss@ufpa.br}{eltonss@ufpa.br}}
\institute[UFPA]{\inst{1} Faculdade de Sistemas de Informação - UFPA Campus Cametá}
\date{\today}

\begin{document}

% ==============================================================================
% 1. TITLE FRAME
% ==============================================================================
\begin{frame}
    \titlepage
    \begin{center}
        \tiny \faGithub\ \href{https://github.com/eltonsarmanho}{github.com/eltonsarmanho} \quad
        \faLinkedinSquare\ \href{https://www.linkedin.com/in/elton-sarmanho-836553185/}{Elton Sarmanho}
    \end{center}
\end{frame}

% ==============================================================================
% 2. ROTEIRO
% ==============================================================================
\begin{frame}{Roteiro}
    \tableofcontents
\end{frame}

% ==============================================================================
\section{Introdução à Classificação}
% ==============================================================================

\begin{frame}{O que é Classificação?}
\begin{columns}[T]
\column{0.5\textwidth}
\begin{block}{Aprendizado Supervisionado}
\begin{itemize}
  \item Aprender a partir de \textbf{exemplos rotulados}
  \item Cada exemplo: atributos + classe
  \item Meta: prever a classe de novos exemplos
\end{itemize}
\end{block}
\vspace{0.4em}
\begin{alertblock}{Definição}
  Dado um conjunto de treino $(\mathbf{x}_i, y_i)$, aprender $f(\mathbf{x}) \rightarrow y$
\end{alertblock}

\column{0.5\textwidth}
\textbf{Aplicações:}
\vspace{0.3em}

\begin{tabular}{cl}
\faEnvelope & \textbf{Spam:} e-mail legítimo ou lixo \\[4pt]
\faHeartbeat & \textbf{Médico:} tumor benigno/maligno \\[4pt]
\faCreditCard & \textbf{Crédito:} aprovação de empréstimo \\[4pt]
\faUsers & \textbf{Churn:} cliente irá cancelar? \\[4pt]
\faShield & \textbf{Fraude:} transação suspeita \\[4pt]
\faCamera & \textbf{Imagem:} reconhecimento de objetos \\
\end{tabular}
\end{columns}
\end{frame}

\begin{frame}{Tipos de Classificação}
\begin{center}
\begin{tabular}{>{\bfseries\color{mineblue}}p{3.2cm}p{4.5cm}p{4.5cm}}
\toprule
Tipo & Descrição & Exemplo \\
\midrule
Binária & 2 classes: $y \in \{0, 1\}$ & Spam vs. Não-spam \\[4pt]
Multiclasse & $k > 2$ classes & Dígito reconhecido (0--9) \\[4pt]
Multirótulo & Múltiplos rótulos simultâneos & Gêneros de um filme \\[4pt]
\bottomrule
\end{tabular}
\end{center}
\vspace{0.5em}
\begin{block}{Nota}
  A maioria dos algoritmos clássicos é projetada para \textbf{classificação binária}. Extensões para multiclasse: \textit{One-vs-Rest} (OvR) e \textit{One-vs-One} (OvO).
\end{block}
\end{frame}

% ==============================================================================
\section{Árvores de Decisão}
% ==============================================================================

\begin{frame}[t]{Árvores de Decisão: Conceito}
\vspace{-0.3em}
\begin{columns}[T]
\column{0.38\textwidth}
\begin{block}{Como funciona?}
\begin{itemize}
  \item \textbf{Nó interno:} teste em um atributo
  \item \textbf{Ramo:} resultado do teste
  \item \textbf{Folha:} classe predita
\end{itemize}
\end{block}
\vspace{0.3em}
Para classificar: percorra da raiz até uma folha seguindo os atributos do exemplo.

\column{0.62\textwidth}
\vspace{-0.3em}
\begin{center}
\scalebox{0.85}{%
\begin{tikzpicture}[
  internal/.style={rectangle,rounded corners,draw=mineblue,fill=mineblue!15,minimum width=2.8cm,minimum height=0.75cm,font=\small,align=center,drop shadow},
  leaf/.style={rectangle,rounded corners,draw=minegreen,fill=minegreen!15,minimum width=2.2cm,minimum height=0.65cm,font=\small\bfseries,align=center,drop shadow},
  edge from parent/.append style={draw,->,thick,mineblue!70},
  level distance=1.7cm,
  level 1/.style={sibling distance=5.5cm},
  level 2/.style={sibling distance=2.8cm}
]
\node[internal] (root) {Renda $>$ 50k?}
  child {
    node[internal] {Histórico\\Limpo?}
    child { node[leaf] {\textcolor{minegreen}{APROVADO}\\(78\%)} edge from parent node[left,font=\tiny\color{minegray}]{Sim} }
    child { node[leaf] {\textcolor{minered}{NEGADO}\\(62\%)} edge from parent node[right,font=\tiny\color{minegray}]{Não} }
    edge from parent node[left,font=\small\color{minegray}]{Não}
  }
  child {
    node[internal] {Emprego\\$>$ 2 anos?}
    child { node[leaf] {\textcolor{minered}{NEGADO}\\(55\%)} edge from parent node[left,font=\tiny\color{minegray}]{Não} }
    child { node[leaf] {\textcolor{minegreen}{APROVADO}\\(91\%)} edge from parent node[right,font=\tiny\color{minegray}]{Sim} }
    edge from parent node[right,font=\small\color{minegray}]{Sim}
  };
\end{tikzpicture}
}
\end{center}
\end{columns}
\end{frame}

\begin{frame}{Algoritmos: ID3, C4.5 e CART}
\begin{center}
\begin{tabular}{>{\bfseries}p{3.2cm}ccc}
\toprule
Característica & \textcolor{mineblue}{ID3} & \textcolor{mineorange}{C4.5} & \textcolor{minegreen}{CART} \\
\midrule
Critério & Entropia/IG & Gain Ratio & Índice Gini \\
Tipo de divisão & Multiway & Multiway & Binária \\
Atrib. contínuos & Não & Sim & Sim \\
Valores ausentes & Não & Sim & Sim \\
Regressão & Não & Não & Sim \\
Poda & Não & Pós-poda & Custo-complexidade \\
\bottomrule
\end{tabular}
\end{center}
\vspace{0.3em}
\begin{columns}[T]
\column{0.5\textwidth}
\begin{block}{Critério de impureza}
  \begin{itemize}
    \item \textbf{Entropia:} $H(S) = -\sum_i p_i \log_2 p_i$
    \item \textbf{Gini:} $G(S) = 1 - \sum_i p_i^2$
  \end{itemize}
\end{block}
\column{0.5\textwidth}
\begin{alertblock}{Poda (Pruning)}
  Evita \textbf{overfitting}: limita profundidade, mínimo de amostras por folha, ou remove ramos após a construção.
\end{alertblock}
\end{columns}
\end{frame}

\begin{frame}{Propriedades das Árvores de Decisão}
\begin{columns}[T]
\column{0.5\textwidth}
\begin{block}{\faThumbsUp\ Vantagens}
\begin{itemize}
  \item \textbf{Interpretabilidade:} regras legíveis por humanos
  \item \textbf{Não-paramétrico:} sem hipóteses sobre a distribuição
  \item \textbf{Sem normalização:} não precisa escalar features
  \item \textbf{Tipos mistos:} numérico e categórico nativamente
\end{itemize}
\end{block}
\column{0.5\textwidth}
\begin{alertblock}{\faWarning\ Limitações}
\begin{itemize}
  \item \textbf{Instabilidade:} pequenas mudanças nos dados $\Rightarrow$ árvore diferente
  \item \textbf{Overfitting:} árvores profundas memorizam o treino
  \item Fronteiras de decisão apenas horizontais/verticais
\end{itemize}
\end{alertblock}
\vspace{0.4em}
\begin{exampleblock}{Solução}
  \textbf{Métodos Ensemble} (Random Forest, Boosting) combinam múltiplas árvores para superar essas limitações.
\end{exampleblock}
\end{columns}
\end{frame}

% ==============================================================================
\section{Métodos Probabilísticos}
% ==============================================================================

\begin{frame}{Naive Bayes}
\begin{columns}[T]
\column{0.48\textwidth}
\begin{block}{Teorema de Bayes}
\vspace{0.3em}
\[\large P(c \mid \mathbf{x}) = \frac{P(\mathbf{x} \mid c)\, P(c)}{P(\mathbf{x})}\]
\vspace{0.1em}
\begin{itemize}
  \item $P(c)$: probabilidade \textit{a priori}
  \item $P(\mathbf{x}|c)$: verossimilhança
  \item $P(c|\mathbf{x})$: probabilidade \textit{a posteriori}
\end{itemize}
\end{block}
\vspace{0.3em}
\textbf{Hipótese ``ingênua''}: atributos são condicionalmente independentes dada a classe.

\column{0.52\textwidth}
\textbf{Variantes:}
\vspace{0.3em}

\begin{tabular}{>{\bfseries\footnotesize}p{2.3cm}p{2.0cm}p{2.0cm}}
\toprule
Variante & Distribuição & Uso \\
\midrule
Gaussiano & Normal & Features contínuas \\[3pt]
Multinomial & Multinomial & Texto (freq. palavras) \\[3pt]
Bernoulli & Bernoulli & Features binárias \\[3pt]
\bottomrule
\end{tabular}
\vspace{0.4em}
\begin{alertblock}{Suavização de Laplace}
  Evita probabilidade zero para valores não vistos no treino: adiciona $\alpha=1$ às contagens.
\end{alertblock}
\end{columns}
\end{frame}

\begin{frame}{Regressão Logística}
\begin{columns}[T]
\column{0.5\textwidth}
\begin{block}{Função Sigmoide}
\vspace{0.3em}
\[\sigma(z) = \frac{1}{1 + e^{-z}}\]
\vspace{0.1em}
onde $z = \mathbf{w}^T \mathbf{x} + b$
\vspace{0.3em}
\begin{itemize}
  \item Saída em $(0, 1)$: probabilidade
  \item Limiar $= 0.5$: classe positiva se $\sigma > 0.5$
\end{itemize}
\end{block}
\vspace{0.3em}
\textbf{Multiclasse:} usa função \textbf{Softmax} para $k$ classes.

\column{0.5\textwidth}
\begin{center}
\begin{tikzpicture}[scale=0.85]
\begin{axis}[
  width=6cm, height=4.5cm,
  xlabel={$z$}, ylabel={$\sigma(z)$},
  xmin=-6, xmax=6, ymin=-0.05, ymax=1.05,
  xtick={-4,-2,0,2,4},
  ytick={0,0.5,1},
  grid=major, grid style={dotted,minegray!40},
  axis lines=left,
  tick label style={font=\tiny},
  label style={font=\small},
]
\addplot[mineblue,very thick,domain=-6:6,samples=100] {1/(1+exp(-x))};
\addplot[dashed,minered,thin] coordinates {(-6,0.5)(6,0.5)};
\end{axis}
\end{tikzpicture}
\end{center}
\begin{exampleblock}{Quando usar?}
  Classificação \textbf{binária linear} com necessidade de probabilidades calibradas.
\end{exampleblock}
\end{columns}
\end{frame}

% ==============================================================================
\section{K-Vizinhos Mais Próximos (KNN)}
% ==============================================================================

\begin{frame}{KNN: K-Vizinhos Mais Próximos}
\begin{columns}[T]
\column{0.5\textwidth}
\begin{block}{Como funciona?}
\begin{enumerate}
  \item Dado novo exemplo $\mathbf{x}$
  \item Calcule a distância para \textbf{todos} os exemplos de treino
  \item Encontre os $k$ mais próximos
  \item \textbf{Vote} pela classe majoritária
\end{enumerate}
\end{block}
\vspace{0.3em}
\textbf{Distância mais comum:} Euclidiana
\[d(\mathbf{a},\mathbf{b}) = \sqrt{\sum_i (a_i - b_i)^2}\]

\column{0.5\textwidth}
\begin{alertblock}{Efeito de $k$}
\begin{itemize}
  \item $k=1$: fronteira irregular \newline \textbf{Overfitting} (muito sensível ao ruído)
  \vspace{0.3em}
  \item $k=n$: fronteira suave \newline \textbf{Underfitting} (ignora estrutura)
  \vspace{0.3em}
  \item $k$ ótimo via \textbf{validação cruzada} \newline (ponto de partida: $k \approx \sqrt{n}$)
\end{itemize}
\end{alertblock}
\begin{exampleblock}{Atenção}
  Requer \textbf{normalização} das features! Distâncias são sensíveis à escala.
\end{exampleblock}
\end{columns}
\end{frame}

% ==============================================================================
\section{Máquinas de Vetores de Suporte (SVM)}
% ==============================================================================

\begin{frame}[t]{SVM: Margem Máxima}
\begin{columns}[T]
\column{0.38\textwidth}
\vspace{0.3em}
\begin{block}{Conceito}
  Encontrar o hiperplano que \textbf{maximiza a margem} entre as classes.
\end{block}
\vspace{0.3em}
\textbf{Vetores de Suporte:} os exemplos mais próximos do hiperplano, sendo os únicos que determinam a fronteira.
\vspace{0.3em}

\begin{alertblock}{Soft Margin}
  Parâmetro $C$ controla o trade-off entre margem e erros de treino.
\end{alertblock}

\column{0.62\textwidth}
\vspace{-0.2em}
\begin{center}
\begin{tikzpicture}[scale=1.0]
% Positive class points (circles - blue)
\foreach \x/\y in {3.5/2.5, 4/3.5, 4.5/2, 5/3, 3.8/1.8}{
  \filldraw[mineblue,fill=mineblue!70] (\x,\y) circle(0.09);
}
% Negative class points (squares - red)
\foreach \x/\y in {1/1.5, 1.5/2.5, 0.8/3, 2/1, 1.8/2.8}{
  \fill[minered,rotate around={45:(\x,\y)}] (\x-0.09,\y-0.09) rectangle (\x+0.09,\y+0.09);
}
% Decision boundary
\draw[very thick,minegreen] (2.5,-0.2) -- (2.5,4.2) node[right,font=\small]{$\mathbf{w}\cdot\mathbf{x}+b=0$};
% Margin boundaries
\draw[dashed,mineblue,thick] (3.2,-0.2) -- (3.2,4.2) node[above,font=\tiny]{$+1$};
\draw[dashed,minered,thick] (1.8,-0.2) -- (1.8,4.2) node[above,font=\tiny]{$-1$};
% Margin annotation
\draw[<->,very thick,mineorange] (1.8,4.0) -- (3.2,4.0) node[midway,above,font=\small]{Margem $\frac{2}{\|\mathbf{w}\|}$};
% Support vectors highlighted
\draw[mineorange,very thick] (3.2,2.5) circle(0.17) node[right,font=\tiny,mineorange]{SV};
\draw[mineorange,very thick] (1.8,2.5) circle(0.17) node[left,font=\tiny,mineorange]{SV};
% Labels
\node[mineblue,font=\small\bfseries] at (4.5,4) {Classe $+1$};
\node[minered,font=\small\bfseries] at (0.8,4) {Classe $-1$};
\end{tikzpicture}
\end{center}
\end{columns}
\end{frame}

\begin{frame}{SVM: Kernel Trick}
\begin{block}{Por que kernels?}
  Quando os dados \textbf{não são linearmente separáveis}, o kernel mapeia para um espaço de maior dimensão \textit{sem calcular o mapeamento explicitamente}.
\end{block}
\vspace{0.5em}
\begin{center}
\begin{tabular}{>{\bfseries\color{mineblue}}p{2.8cm}p{4.5cm}p{4cm}}
\toprule
Kernel & Fórmula & Quando usar \\
\midrule
Linear & $K(\mathbf{x},\mathbf{z}) = \mathbf{x} \cdot \mathbf{z}$ & Dados linearmente separáveis, texto \\[4pt]
Polinomial & $K(\mathbf{x},\mathbf{z}) = (\mathbf{x} \cdot \mathbf{z} + c)^d$ & Interações de grau $d$ \\[4pt]
RBF (Gaussiano) & $K(\mathbf{x},\mathbf{z}) = e^{-\gamma\|\mathbf{x}-\mathbf{z}\|^2}$ & Fronteiras complexas (padrão) \\[4pt]
Sigmoide & $K(\mathbf{x},\mathbf{z}) = \tanh(\kappa\mathbf{x}\cdot\mathbf{z} + c)$ & Inspirado em redes neurais \\[4pt]
\bottomrule
\end{tabular}
\end{center}
\vspace{0.3em}
\begin{exampleblock}{Dica prática}
  Comece com \textbf{RBF} (default no scikit-learn). Use \textbf{Linear} para texto esparso ou dados de alta dimensão.
\end{exampleblock}
\end{frame}

% ==============================================================================
\section{Métodos Ensemble}
% ==============================================================================

\begin{frame}[t]{Métodos Ensemble: Visão Geral}
\begin{columns}[T]
\column{0.38\textwidth}
\vspace{0.3em}
\begin{block}{Princípio}
  \textbf{Muitos modelos} simples $>$ \textbf{um modelo} complexo
\end{block}
\vspace{0.3em}
\begin{itemize}
  \item \textbf{Bagging:} paralelo, reduz \textit{variância}
  \vspace{0.2em}
  \item \textbf{Boosting:} sequencial, reduz \textit{viés}
\end{itemize}

\column{0.62\textwidth}
\vspace{-0.2em}
\begin{center}
\scalebox{0.75}{%
\begin{tikzpicture}[scale=0.95]
% BAGGING side
\node[font=\bfseries\small,color=mineblue] at (-3.2,4.0) {BAGGING};
\node[draw=minegray,rounded corners,font=\tiny\bfseries,minimum width=1.8cm,minimum height=0.5cm,fill=minegray!10] (dbag) at (-3.2,3.3) {Dataset $\mathcal{D}$};

\node[draw=mineblue,fill=mineblue!10,rounded corners,font=\tiny,minimum width=1.4cm,minimum height=0.5cm] (b1) at (-4.5,2.2) {Bootstrap 1};
\node[draw=mineblue,fill=mineblue!10,rounded corners,font=\tiny,minimum width=1.4cm,minimum height=0.5cm] (b2) at (-3.2,2.2) {Bootstrap 2};
\node[draw=mineblue,fill=mineblue!10,rounded corners,font=\tiny,minimum width=1.4cm,minimum height=0.5cm] (b3) at (-1.9,2.2) {Bootstrap 3};

\draw[->,minegray,thin] (dbag) -- (b1); \draw[->,minegray,thin] (dbag) -- (b2); \draw[->,minegray,thin] (dbag) -- (b3);

\node[draw=mineblue,fill=mineblue!20,rounded corners,font=\tiny,minimum width=1.4cm,minimum height=0.5cm] (bt1) at (-4.5,1.2) {Modelo 1};
\node[draw=mineblue,fill=mineblue!20,rounded corners,font=\tiny,minimum width=1.4cm,minimum height=0.5cm] (bt2) at (-3.2,1.2) {Modelo 2};
\node[draw=mineblue,fill=mineblue!20,rounded corners,font=\tiny,minimum width=1.4cm,minimum height=0.5cm] (bt3) at (-1.9,1.2) {Modelo 3};

\draw[->,mineblue,thin] (b1)--(bt1); \draw[->,mineblue,thin] (b2)--(bt2); \draw[->,mineblue,thin] (b3)--(bt3);

\node[draw=minegreen,fill=minegreen!10,rounded corners,font=\tiny,minimum width=2.2cm,minimum height=0.5cm] (bagg) at (-3.2,0.2) {Votação / Média};
\draw[->,minegreen,thin] (bt1)--(bagg); \draw[->,minegreen,thin] (bt2)--(bagg); \draw[->,minegreen,thin] (bt3)--(bagg);

% Divider
\draw[dashed,minegray!60,thick] (0,4.5) -- (0,-0.3);

% BOOSTING side
\node[font=\bfseries\small,color=mineorange] at (3.2,4.0) {BOOSTING};
\node[draw=minegray,rounded corners,font=\tiny\bfseries,minimum width=1.8cm,minimum height=0.5cm,fill=minegray!10] (dbst) at (3.2,3.3) {Dataset $\mathcal{D}$};

\node[draw=mineorange,fill=mineorange!10,rounded corners,font=\tiny,minimum width=2cm,minimum height=0.5cm] (bst1) at (3.2,2.5) {Modelo 1};
\draw[->,minegray,thin] (dbst)--(bst1);

\node[draw=mineorange,fill=mineorange!15,rounded corners,font=\tiny,minimum width=2cm,minimum height=0.5cm,align=center] (bst2) at (3.2,1.7) {Modelo 2\\{\color{minegray}(foco erros M1)}};
\draw[->,mineorange,thin] (bst1)--(bst2);

\node[draw=mineorange,fill=mineorange!20,rounded corners,font=\tiny,minimum width=2cm,minimum height=0.5cm,align=center] (bst3) at (3.2,0.8) {Modelo 3\\{\color{minegray}(foco erros M2)}};
\draw[->,mineorange,thin] (bst2)--(bst3);

\node[draw=minegreen,fill=minegreen!10,rounded corners,font=\tiny,minimum width=2.2cm,minimum height=0.5cm,align=center] (bstf) at (3.2,0.0) {Combinação\\ponderada};
\draw[->,minegreen,thin] (bst3)--(bstf);
\end{tikzpicture}
}
\end{center}
\end{columns}
\end{frame}

\begin{frame}{Random Forest}
\begin{columns}[T]
\column{0.5\textwidth}
\begin{block}{Como funciona?}
\begin{enumerate}
  \item \textbf{Bootstrap:} criar $T$ amostras com reposição
  \item \textbf{Árvore aleatória:} em cada divisão, avaliar apenas $m \approx \sqrt{p}$ features aleatórias
  \item \textbf{Votação:} predição final = classe majoritária das $T$ árvores
\end{enumerate}
\end{block}
\vspace{0.4em}
A aleatoriedade nas features \textbf{descorrelaciona} as árvores, reduzindo a variância.

\column{0.5\textwidth}
\begin{exampleblock}{Vantagens}
\begin{itemize}
  \item Muito robusto a overfitting
  \item Lida bem com muitas features
  \item \textbf{Importância de features} (MDI, MDA)
  \item Poucas decisões de hiperparâmetros
\end{itemize}
\end{exampleblock}
\vspace{0.4em}
\begin{alertblock}{Limitação}
  Menos interpretável que uma única árvore. Para explicabilidade: usar \textbf{SHAP}.
\end{alertblock}
\end{columns}
\end{frame}

\begin{frame}{XGBoost, LightGBM e CatBoost}
\begin{block}{Gradient Boosted Decision Trees (GBDT)}
  Constroem árvores \textbf{sequencialmente}, cada uma corrigindo os erros da anterior. Dominam benchmarks em dados tabulares.
\end{block}
\vspace{0.4em}
\begin{center}
\begin{tabular}{>{\bfseries\color{mineblue}}p{2.8cm}p{4cm}p{4.5cm}}
\toprule
Algoritmo & Inovação chave & Destaque \\
\midrule
XGBoost & Regularização L1/L2 & Competições Kaggle desde 2015 \\[3pt]
LightGBM & GOSS + crescimento leaf-wise & $\sim$10$\times$ mais rápido em dados grandes \\[3pt]
CatBoost & Ordered target encoding & Suporte nativo a categóricas \\[3pt]
\bottomrule
\end{tabular}
\end{center}
\vspace{0.4em}
\begin{exampleblock}{Consulte o Material Didático e o Notebook Colab}
  Implementações completas com tuning de hiperparâmetros e comparação de desempenho estão disponíveis no \textbf{Módulo 3} do material didático.
\end{exampleblock}
\end{frame}

% ==============================================================================
\section{Avaliação de Modelos}
% ==============================================================================

\begin{frame}[t]{Matriz de Confusão}
\begin{columns}[T]
\column{0.45\textwidth}
\vspace{0.3em}
\begin{block}{O que é?}
  Tabela que organiza as predições vs. rótulos reais. Revela \textbf{tipos de erros} cometidos pelo modelo.
\end{block}
\vspace{0.3em}
\begin{itemize}
  \item \textbf{VP}: acertou positivo
  \item \textbf{VN}: acertou negativo
  \item \textbf{FP}: disse positivo, era negativo
  \item \textbf{FN}: disse negativo, era positivo
\end{itemize}

\column{0.55\textwidth}
\begin{center}
\scalebox{0.55}{%
\begin{tikzpicture}
  \def\cw{3.2}
  \def\hh{0.85}
  \def\ch{1.6}

  % Preenchimento
  \fill[minegray!12]  (0,0)        rectangle (\cw,-\hh);
  \fill[mineblue!20]  (\cw,0)      rectangle ({2*\cw},-\hh);
  \fill[minered!15]   ({2*\cw},0)  rectangle ({3*\cw},-\hh);
  \fill[mineblue!20]  (0,-\hh)              rectangle (\cw,{-\hh-\ch});
  \fill[minegreen!25] (\cw,-\hh)            rectangle ({2*\cw},{-\hh-\ch});
  \fill[minered!15]   ({2*\cw},-\hh)        rectangle ({3*\cw},{-\hh-\ch});
  \fill[minered!15]   (0,{-\hh-\ch})        rectangle (\cw,{-\hh-2*\ch});
  \fill[minered!15]   (\cw,{-\hh-\ch})      rectangle ({2*\cw},{-\hh-2*\ch});
  \fill[minegreen!25] ({2*\cw},{-\hh-\ch})  rectangle ({3*\cw},{-\hh-2*\ch});

  % Grade
  \foreach \x in {0,\cw,{2*\cw},{3*\cw}}{
    \draw[minegray!60,line width=0.5pt] (\x,0) -- (\x,{-\hh-2*\ch});
  }
  \foreach \y in {0,-\hh,{-\hh-\ch},{-\hh-2*\ch}}{
    \draw[minegray!60,line width=0.5pt] (0,\y) -- ({3*\cw},\y);
  }

  % Cabeçalho
  \node[font=\small\bfseries] at ({0.5*\cw},{-0.5*\hh}) {};
  \node[font=\small\bfseries,text=mineblue!80!black]
    at ({1.5*\cw},{-0.5*\hh}) {Previsto $+$};
  \node[font=\small\bfseries,text=minered!80!black]
    at ({2.5*\cw},{-0.5*\hh}) {Previsto $-$};

  % Linha Real +
  \node[font=\small\bfseries,text=mineblue!80!black]
    at ({0.5*\cw},{-\hh-0.5*\ch}) {Real $+$};
  \node[font=\small,align=center]
    at ({1.5*\cw},{-\hh-0.5*\ch})
    {\textbf{\large VP}\\[2pt]\scriptsize(Verdadeiro Positivo)};
  \node[font=\small,align=center]
    at ({2.5*\cw},{-\hh-0.5*\ch})
    {\textbf{\large FN}\\[2pt]\scriptsize(Falso Negativo)};

  % Linha Real -
  \node[font=\small\bfseries,text=minered!80!black]
    at ({0.5*\cw},{-\hh-1.5*\ch}) {Real $-$};
  \node[font=\small,align=center]
    at ({1.5*\cw},{-\hh-1.5*\ch})
    {\textbf{\large FP}\\[2pt]\scriptsize(Falso Positivo)};
  \node[font=\small,align=center]
    at ({2.5*\cw},{-\hh-1.5*\ch})
    {\textbf{\large VN}\\[2pt]\scriptsize(Verdadeiro Negativo)};
\end{tikzpicture}
}
\end{center}
\end{columns}
\end{frame}

\begin{frame}{Métricas de Avaliação}
\begin{columns}[T]
\column{0.5\textwidth}
\begin{block}{Fórmulas}
\begin{align*}
\text{Acurácia} &= \frac{VP + VN}{VP+VN+FP+FN} \\[6pt]
\text{Precisão} &= \frac{VP}{VP + FP} \\[6pt]
\text{Recall} &= \frac{VP}{VP + FN} \\[6pt]
\text{F1-Score} &= 2 \cdot \frac{\text{P} \cdot \text{R}}{\text{P} + \text{R}}
\end{align*}
\end{block}

\column{0.5\textwidth}
\begin{exampleblock}{Exemplo: Detecção de Doença}
Imagine 100 pacientes (10 doentes, 90 saudáveis):
\begin{itemize}
  \item Modelo detectou 8 dos 10 doentes (VP=8)
  \item 5 saudáveis foram ditos doentes (FP=5)
  \item \textbf{Precisão} = 8/13 = 0,62
  \item \textbf{Recall} = 8/10 = 0,80
  \item \textbf{F1} = 0,70
\end{itemize}
\end{exampleblock}
\begin{alertblock}{Acurácia engana!}
  Se 95\% dos e-mails são legítimos, um modelo que classifica \textit{tudo} como legítimo tem 95\% de acurácia.
\end{alertblock}
\end{columns}
\end{frame}

\begin{frame}{Quando Usar Cada Métrica?}
\footnotesize
\begin{center}
\begin{tabular}{>{\bfseries\color{mineblue}}p{3cm}p{3.5cm}p{5cm}}
\toprule
Situação & Métrica & Motivo \\
\midrule
Classes balanceadas & Acurácia & Simples e intuitiva \\[4pt]
FP tem alto custo & Precisão & Spam: não bloquear e-mail legítimo \\[4pt]
FN tem alto custo & Recall & Diagnóstico: não perder paciente doente \\[4pt]
Equilíbrio P/R & F1-Score & Geral para dados desbalanceados \\[4pt]
Threshold incerto & AUC-ROC & Avalia ranking de probabilidades \\[4pt]
Alta desbalanceamento & AUC-PR & Fraude, doenças raras \\[4pt]
\bottomrule
\end{tabular}
\end{center}

\end{frame}

% --- Frame: Curva de Aprendizado ---
\begin{frame}[t]{Curva de Aprendizado: O que o Modelo Aprende?}
\footnotesize
\begin{columns}[T]
\column{0.60\textwidth}
\vspace{0.3em}
\begin{block}{O que é?}
  Plota o desempenho (acurácia) em função do \textbf{tamanho do conjunto de treino}.\\[4pt]
  Avalia se o modelo está \textbf{aprendendo de fato} ou apenas memorizando dados.
\end{block}
\vspace{0.4em}
\begin{itemize}\small
  \item \textbf{Curva de treino:} desempenho nos dados de treino
  \item \textbf{Curva de validação:} desempenho em dados novos
\end{itemize}
\vspace{0.4em}
\begin{alertblock}{Diagnóstico pelo gap}
  O \textbf{gap} entre as duas curvas revela se há overfitting, underfitting ou aprendizado saudável.
\end{alertblock}
\column{0.60\textwidth}
\vspace{-0.2em}
\begin{center}
\scalebox{0.82}{%
\begin{tikzpicture}
\begin{axis}[
  width=8cm, height=5.5cm,
  xlabel={Tamanho do Treino},
  ylabel={Acurácia},
  xmin=50,xmax=500,ymin=0.50,ymax=1.02,
  xtick={100,200,300,400,500},
  ytick={0.6,0.7,0.8,0.9,1.0},
  yticklabels={60\%,70\%,80\%,90\%,100\%},
  grid=major, grid style={dotted,minegray!40},
  axis lines=left,
  tick label style={font=\tiny},
  label style={font=\small},
  legend style={font=\tiny,at={(0.98,0.32)},anchor=east},
]
\addplot[minered,very thick,smooth] coordinates {
  (50,0.99)(100,0.97)(150,0.95)(200,0.93)(300,0.91)(400,0.90)(500,0.89)};
\addlegendentry{Treino (overfit)}
\addplot[minered,dashed,thick,smooth] coordinates {
  (50,0.60)(100,0.65)(150,0.68)(200,0.70)(300,0.71)(400,0.72)(500,0.73)};
\addlegendentry{Validação (overfit)}
\addplot[mineblue,very thick,smooth] coordinates {
  (50,0.95)(100,0.91)(150,0.89)(200,0.87)(300,0.86)(400,0.85)(500,0.85)};
\addlegendentry{Treino (bom)}
\addplot[mineblue,dashed,thick,smooth] coordinates {
  (50,0.65)(100,0.74)(150,0.79)(200,0.82)(300,0.84)(400,0.85)(500,0.85)};
\addlegendentry{Validação (bom)}
\draw[<->,mineorange,thick] (axis cs:500,0.89) -- (axis cs:500,0.73)
  node[midway,right,font=\tiny,mineorange]{gap grande};
\end{axis}
\end{tikzpicture}
}
\end{center}
\end{columns}
\end{frame}

% --- Frame: Diagnóstico da Curva de Aprendizado ---
\begin{frame}[shrink=5]{Curva de Aprendizado: Três Cenários}
\begin{center}
\scalebox{0.72}{%
\begin{tikzpicture}

% ---- OVERFITTING ----
\node[font=\bfseries\small,minered] at (2.2,5.0) {Overfitting};
\node[font=\tiny,minered] at (2.2,4.6) {(Alta Variância)};
\begin{axis}[
  at={(0cm,0cm)},
  width=5.5cm,height=4.5cm,
  xmin=0,xmax=10,ymin=0.40,ymax=1.05,
  xtick=\empty,
  ytick={0.5,0.75,1.0},yticklabels={50\%,75\%,100\%},
  tick label style={font=\tiny},
  axis lines=left,
  grid=major,grid style={dotted,minegray!30},
]
\addplot[minered,very thick,smooth] coordinates
  {(1,0.99)(3,0.97)(5,0.95)(7,0.93)(10,0.92)};
\addplot[minered,dashed,thick,smooth] coordinates
  {(1,0.55)(3,0.60)(5,0.63)(7,0.65)(10,0.66)};
\draw[<->,mineorange,thin] (axis cs:8,0.92) -- (axis cs:8,0.66)
  node[midway,right,font=\tiny]{gap};
\end{axis}
\node[font=\tiny,align=center,minered,text width=4.5cm] at (2.2,-0.8)
  {Treino ótimo, validação baixa\\ \textbf{Solução:} mais dados, regularização,\\ menos features};

% ---- UNDERFITTING ----
\node[font=\bfseries\small,mineorange] at (9.7,5.0) {Underfitting};
\node[font=\tiny,mineorange] at (9.7,4.6) {(Alto Viés)};
\begin{axis}[
  at={(7.5cm,0cm)},
  width=5.5cm,height=4.5cm,
  xmin=0,xmax=10,ymin=0.40,ymax=1.05,
  xtick=\empty,
  ytick={0.5,0.75,1.0},yticklabels=\empty,
  tick label style={font=\tiny},
  axis lines=left,
  grid=major,grid style={dotted,minegray!30},
]
\addplot[mineorange,very thick,smooth] coordinates
  {(1,0.72)(3,0.70)(5,0.69)(7,0.68)(10,0.67)};
\addplot[mineorange,dashed,thick,smooth] coordinates
  {(1,0.65)(3,0.66)(5,0.66)(7,0.67)(10,0.67)};
\end{axis}
\node[font=\tiny,align=center,mineorange,text width=4.5cm] at (9.7,-0.8)
  {Ambas as curvas baixas, gap pequeno\\ \textbf{Solução:} modelo mais complexo,\\ mais features, menos regularização};

% ---- BOM MODELO ----
\node[font=\bfseries\small,minegreen] at (17.2,5.0) {Modelo Ideal};
\node[font=\tiny,minegreen] at (17.2,4.6) {(Balanceado)};
\begin{axis}[
  at={(15cm,0cm)},
  width=5.5cm,height=4.5cm,
  xmin=0,xmax=10,ymin=0.40,ymax=1.05,
  xtick=\empty,
  ytick={0.5,0.75,1.0},yticklabels=\empty,
  tick label style={font=\tiny},
  axis lines=left,
  grid=major,grid style={dotted,minegray!30},
]
\addplot[minegreen,very thick,smooth] coordinates
  {(1,0.90)(3,0.87)(5,0.86)(7,0.85)(10,0.85)};
\addplot[minegreen,dashed,thick,smooth] coordinates
  {(1,0.65)(3,0.75)(5,0.81)(7,0.84)(10,0.85)};
\end{axis}
\node[font=\tiny,align=center,minegreen,text width=4.5cm] at (17.2,-0.8)
  {Curvas convergem em valor alto\\ Gap pequeno ao final\\ \textbf{Generaliza bem!}};

% Legenda
\draw[black,very thick] (6.5,3.8) -- (7.2,3.8) node[right,font=\tiny]{Treino};
\draw[black,dashed,thick] (6.5,3.3) -- (7.2,3.3) node[right,font=\tiny]{Validação};
\end{tikzpicture}
}
\end{center}
\end{frame}

% --- Frame: Curva ROC ---
\begin{frame}[t]{}
\vspace{-0.83cm}
\begin{columns}[T]
\column{0.42\textwidth}
\vspace{0.2em}
\begin{block}{O que é a Curva ROC?}
  Plota a \textbf{Taxa de Verdadeiros Positivos} (Recall/Sensibilidade) contra a \textbf{Taxa de Falsos Positivos} para \textit{todos os possíveis limiares} de decisão.
\end{block}
\vspace{0.2em}
\begin{itemize}\small
  \item $TVP = \dfrac{VP}{VP+FN}$ \quad (eixo Y)
  \vspace{0.2em}
  \item $TFP = \dfrac{FP}{FP+VN}$ \quad (eixo X)
\end{itemize}
\vspace{0.3em}
\begin{block}{AUC: Área Sob a Curva}
  \begin{itemize}\small
    \item AUC $= 1.0$: modelo \textbf{perfeito}
    \item AUC $= 0.5$: modelo \textbf{aleatório}
    \item AUC $< 0.5$: pior que o acaso
  \end{itemize}
\end{block}

\column{0.58\textwidth}
\begin{center}
\scalebox{0.85}{%
\begin{tikzpicture}
\begin{axis}[
  width=7.5cm, height=6.5cm,
  xlabel={Taxa de Falsos Positivos (TFP)},
  ylabel={Taxa de Verdadeiros Positivos (TVP)},
  xmin=0,xmax=1,ymin=0,ymax=1.05,
  xtick={0,0.2,0.4,0.6,0.8,1.0},
  ytick={0,0.2,0.4,0.6,0.8,1.0},
  xticklabels={0,20\%,40\%,60\%,80\%,100\%},
  yticklabels={0,20\%,40\%,60\%,80\%,100\%},
  grid=major, grid style={dotted,minegray!40},
  axis lines=left,
  tick label style={font=\tiny},
  label style={font=\small},
  legend style={font=\tiny,at={(0.98,0.20)},anchor=east},
]
\addplot[minegray,dashed,thick] coordinates {(0,0)(1,1)};
\addlegendentry{Aleatório (AUC=0.50)}
\addplot[mineorange,thick,smooth] coordinates {
  (0,0)(0.05,0.40)(0.10,0.60)(0.20,0.72)(0.30,0.80)(0.50,0.88)(0.70,0.93)(1.0,1.0)};
\addlegendentry{Modelo B (AUC=0.82)}
\addplot[mineblue,very thick,smooth] coordinates {
  (0,0)(0.02,0.60)(0.05,0.80)(0.10,0.88)(0.20,0.93)(0.40,0.97)(0.70,0.99)(1.0,1.0)};
\addlegendentry{Modelo A (AUC=0.95)}
\addplot[minegreen,thick] coordinates {(0,0)(0,1)(1,1)};
\addlegendentry{Perfeito (AUC=1.00)}
\filldraw[mineblue] (axis cs:0.10,0.88) circle(2.5pt)
  node[above right,font=\tiny,mineblue]{limiar escolhido};
\end{axis}
\end{tikzpicture}
}
\end{center}
\end{columns}
\end{frame}

% --- Frame: AUC-ROC Exemplo Prático ---
\begin{frame}{}
  \vspace{-0.75cm}
\begin{block}{Cenário}
  \footnotesize
  Classificador atribui probabilidade de spam a 10 e-mails (5 spam, 5 legítimos). Variando o limiar, calculamos os pontos da curva ROC.
\end{block}
\vspace{0.3em}
\begin{columns}[T]
\column{0.50\textwidth}
\begin{center}\footnotesize
\renewcommand{\arraystretch}{1.3}
\begin{tabular}{ccccc}
\toprule
\textbf{Limiar} & \textbf{VP} & \textbf{FP} & \textbf{TVP} & \textbf{TFP} \\
\midrule
0.9 & 2 & 0 & 0.40 & 0.00 \\
0.7 & 3 & 1 & 0.60 & 0.20 \\
\rowcolor{mineblue!12}
0.5 & 4 & 1 & 0.80 & 0.20 \\
0.3 & 5 & 2 & 1.00 & 0.40 \\
0.1 & 5 & 5 & 1.00 & 1.00 \\
\bottomrule
\end{tabular}
\end{center}
\vspace{0.2em}
\begin{alertblock}{}
  \footnotesize
  \textbf{Limiar baixo:} captura mais spam (TVP$\uparrow$), mas bloqueia mais legítimos (TFP$\uparrow$).\\[3pt]
  O ponto azul na tabela (limiar 0.5) equilibra os erros.
\end{alertblock}
\column{0.50\textwidth}
\begin{block}{AUC-ROC vs AUC-PR}
  \begin{itemize}\footnotesize
    \item \textbf{AUC-ROC:} ideal para classes \textbf{balanceadas}; avalia separação geral entre classes
    \item \textbf{AUC-PR} (Precisão $\times$ Recall): preferível em dados \textbf{muito desbalanceados} (fraude, doenças raras)
  \end{itemize}
\end{block}
\vspace{0.3em}
\begin{exampleblock}{Referência rápida}
  \footnotesize
  AUC $\geq 0.90$: excelente\\
  AUC $\geq 0.80$: bom\\
  AUC $\geq 0.70$: aceitável\\
  AUC $< 0.60$: rever o modelo
\end{exampleblock}
\end{columns}
\end{frame}

% --- Frame: Cálculo detalhado passo-a-passo ---
\begin{frame}[shrink=5]{ROC Spam: Passo a Passo — Como Calcular TVP e TFP}
\begin{block}{10 e-mails classificados (5 spam \textcolor{minered}{$\bullet$}, 5 legítimos \textcolor{mineblue}{$\circ$})}
  \footnotesize
  O modelo atribui uma probabilidade de spam a cada e-mail. Ordenamos do maior ao menor:
  \begin{center}
  \renewcommand{\arraystretch}{1.15}
  \begin{tabular}{cccccc}
  \toprule
  \textbf{E-mail} & \textbf{P(spam)} & \textbf{Rótulo Real} & \textbf{VP acum.} & \textbf{FP acum.} & \textbf{TVP / TFP} \\
  \midrule
  $e_1$ & 0.95 & \textcolor{minered}{\textbf{Spam}}   & 1 & 0 & 0.20 / 0.00 \\
  $e_2$ & 0.88 & \textcolor{minered}{\textbf{Spam}}   & 2 & 0 & 0.40 / 0.00 \\
  $e_3$ & 0.80 & \textcolor{mineblue}{Legítimo}        & 2 & 1 & 0.40 / 0.20 \\
  $e_4$ & 0.75 & \textcolor{minered}{\textbf{Spam}}   & 3 & 1 & 0.60 / 0.20 \\
  $e_5$ & 0.65 & \textcolor{minered}{\textbf{Spam}}   & 4 & 1 & 0.80 / 0.20 \\
  $e_6$ & 0.55 & \textcolor{mineblue}{Legítimo}        & 4 & 2 & 0.80 / 0.40 \\
  $e_7$ & 0.40 & \textcolor{minered}{\textbf{Spam}}   & 5 & 2 & 1.00 / 0.40 \\
  $e_8$ & 0.30 & \textcolor{mineblue}{Legítimo}        & 5 & 3 & 1.00 / 0.60 \\
  $e_9$ & 0.20 & \textcolor{mineblue}{Legítimo}        & 5 & 4 & 1.00 / 0.80 \\
  $e_{10}$ & 0.10 & \textcolor{mineblue}{Legítimo}     & 5 & 5 & 1.00 / 1.00 \\
  \bottomrule
  \end{tabular}
  \end{center}
\end{block}
\vspace{0.2em}
\begin{columns}[T]
\column{0.50\textwidth}
\begin{exampleblock}{Fórmula aplicada}
  \footnotesize
  $TVP = \dfrac{\text{VP acumulado}}{5 \text{ spams}} \qquad TFP = \dfrac{\text{FP acumulado}}{5 \text{ legítimos}}$\\[4pt]
  Cada linha = um ponto na curva ROC ao usar aquela prob.\ como limiar.
\end{exampleblock}
\column{0.50\textwidth}
\begin{alertblock}{Leitura da tabela}
  \footnotesize
  \textbf{Limiar 0.65} (linha $e_5$): captura 4 dos 5 spams (TVP=80\%) errando apenas 1 legítimo (TFP=20\%). Ótimo ponto de operação!
\end{alertblock}
\end{columns}
\end{frame}

% --- Frame: Curva ROC plotada do exemplo ---
\begin{frame}{ROC Spam: Curva Plotada com os 10 Pontos}
  \vspace{-0.8cm}
\begin{columns}[T]
\column{0.45\textwidth}
\vspace{0.3em}
\begin{block}{Pontos da Curva}
  \footnotesize
  \renewcommand{\arraystretch}{1.2}
  \begin{tabular}{cccc}
  \toprule
  \textbf{Pt} & \textbf{Limiar} & \textbf{TFP} & \textbf{TVP} \\
  \midrule
  $p_0$ & — & 0.00 & 0.00 \\
  $p_1$ & 0.95 & 0.00 & 0.20 \\
  $p_2$ & 0.88 & 0.00 & 0.40 \\
  $p_3$ & 0.80 & 0.20 & 0.40 \\
  $p_4$ & 0.75 & 0.20 & 0.60 \\
  \rowcolor{minegreen!15}
  $p_5$ & 0.65 & 0.20 & 0.80 \\
  $p_6$ & 0.55 & 0.40 & 0.80 \\
  $p_7$ & 0.40 & 0.40 & 1.00 \\
  $p_8$ & 0.30 & 0.60 & 1.00 \\
  $p_9$ & 0.20 & 0.80 & 1.00 \\
  $p_{10}$ & 0.10 & 1.00 & 1.00 \\
  \bottomrule
  \end{tabular}
\end{block}

\column{0.55\textwidth}
\begin{center}
\begin{tikzpicture}
\begin{axis}[
  width=7.2cm, height=6.8cm,
  xlabel={Taxa de Falsos Positivos (TFP)},
  ylabel={Taxa de Verdadeiros Positivos (TVP)},
  xmin=-0.02,xmax=1.05,ymin=-0.02,ymax=1.08,
  xtick={0,0.2,0.4,0.6,0.8,1.0},
  ytick={0,0.2,0.4,0.6,0.8,1.0},
  xticklabels={0\%,20\%,40\%,60\%,80\%,100\%},
  yticklabels={0\%,20\%,40\%,60\%,80\%,100\%},
  grid=major,grid style={dotted,minegray!40},
  axis lines=left,
  tick label style={font=\tiny},
  label style={font=\small},
]
% Diagonal aleatório
\addplot[minegray,dashed,thick] coordinates {(0,0)(1,1)};
% Área sob a curva (preenchida)
\addplot[mineblue!20,fill=mineblue!15] coordinates {
  (0,0)(0,0.2)(0,0.4)(0.2,0.4)(0.2,0.6)(0.2,0.8)(0.4,0.8)(0.4,1.0)(0.6,1.0)(0.8,1.0)(1.0,1.0)(1.0,0)(0,0)
} -- cycle;
% Curva ROC
\addplot[mineblue,very thick] coordinates {
  (0,0)(0,0.2)(0,0.4)(0.2,0.4)(0.2,0.6)(0.2,0.8)(0.4,0.8)(0.4,1.0)(0.6,1.0)(0.8,1.0)(1.0,1.0)
};
% Pontos
\addplot[mineblue,only marks,mark=*,mark size=2.5pt] coordinates {
  (0,0)(0,0.2)(0,0.4)(0.2,0.4)(0.2,0.6)(0.2,0.8)(0.4,0.8)(0.4,1.0)(0.6,1.0)(0.8,1.0)(1.0,1.0)
};
% Ponto ideal destacado (limiar 0.65, p5)
\filldraw[minegreen,thick] (axis cs:0.2,0.8) circle(4pt);
\node[above right,font=\tiny,minegreen] at (axis cs:0.2,0.8) {$p_5$ (limiar 0.65)};
% AUC label
\node[font=\small\bfseries,mineblue] at (axis cs:0.65,0.35) {AUC $\approx 0.88$};
% Labels nos eixos extremos
\node[font=\tiny,minegray] at (axis cs:0.5,0.45) {aleatório};
\end{axis}
\end{tikzpicture}
\end{center}
\end{columns}
\end{frame}

% --- Frame: Por que usar AUC? ---
\begin{frame}{}
  \vspace{-0.8cm}
  \footnotesize
\begin{columns}[T]
\column{0.50\textwidth}

\begin{block}{\faThumbsUp\ Benefícios da AUC-ROC}
  \begin{itemize}
    \item \textbf{Independente do limiar:} avalia o modelo para \textit{todos} os limiares de decisão simultaneamente, sem fixar um valor de corte
    \vspace{0.2em}
    \item \textbf{Invariante à escala:} não depende dos valores absolutos das probabilidades preditas, apenas da ordenação
    \vspace{0.2em}
    \item \textbf{Robusta a classes balanceadas:} compara modelos sem viés causado pela proporção entre classes
    \vspace{0.2em}
    \item \textbf{Comparação direta:} um número único (0 a 1) para comparar modelos diferentes de forma objetiva
    \vspace{0.2em}
    \item \textbf{Interpretação probabilística:} AUC = probabilidade de o modelo ranquear um positivo aleatório acima de um negativo aleatório
  \end{itemize}
\end{block}

\column{0.50\textwidth}

\begin{exampleblock}{Interpretação Probabilística (resumida)}
  \footnotesize
  $\text{AUC} = P(\hat{p}_{\text{spam}} > \hat{p}_{\text{legítimo}})$\\[6pt]
  \textbf{AUC = 0.88} significa: em 88\% dos pares (spam, legítimo) sorteados aleatoriamente, o modelo atribui probabilidade \textit{maior} ao spam, mesmo sem definir nenhum limiar.
\end{exampleblock}

\vspace{0.4em}

\begin{alertblock}{\faWarning\ Limitações da AUC-ROC}
  \begin{itemize}\footnotesize
    \item \textbf{Classes muito desbalanceadas:} muitos negativos inflam TFP e mascaram desempenho real → prefira \textbf{AUC-PR}
    \item \textbf{Não informa o limiar ideal:} é necessária análise adicional (curva, custo de erro) para escolher o ponto de operação
    \item \textbf{Não é suficiente sozinha:} AUC alta não garante bom desempenho em um limiar específico relevante ao negócio
  \end{itemize}
\end{alertblock}

\vspace{0.3em}
\begin{center}
\small
\begin{tabular}{lcc}
\toprule
\textbf{Situação} & \textbf{Use} & \textbf{Por quê} \\
\midrule
Classes balanceadas & AUC-ROC & Avalia ranqueamento geral \\
Classes desbalanceadas & AUC-PR & Foca na classe minoritária \\
Custo assimétrico & Curva ROC & Escolhe limiar pelo custo \\
\bottomrule
\end{tabular}
\end{center}

\end{columns}
\end{frame}

\begin{frame}{Validação Cruzada (k-fold)}
  \footnotesize

\begin{columns}[T]
\column{0.55\textwidth}
\begin{block}{Como funciona? (k=5)}
\begin{enumerate}
  \item Divide os dados em $k$ partes (\textit{folds})
  \item Em cada rodada: usa $k-1$ folds para treino, 1 para validação
  \item Repete $k$ vezes (cada fold serve como validação uma vez)
  \item Resultado: média e desvio-padrão das $k$ avaliações
\end{enumerate}
\end{block}
\vspace{0.3em}
\textbf{Variante:} \textit{Stratified k-fold} mantém proporção de classes em cada fold, essencial para dados desbalanceados.

\column{0.45\textwidth}
\begin{center}
\begin{tikzpicture}[scale=0.85]
\foreach \i in {1,...,5}{
  \foreach \j in {1,...,5}{
    \pgfmathsetmacro{\isequal}{int(\i==\j)}
    \ifnum\isequal=1
      \fill[mineorange!50] ({(\j-1)*1.1},{\i*0.6}) rectangle ({(\j)*1.1-0.05},{(\i)*0.6+0.5});
      \node[font=\tiny,white] at ({(\j-0.5)*1.1-0.025},{(\i)*0.6+0.25}) {Val};
    \else
      \fill[mineblue!30] ({(\j-1)*1.1},{\i*0.6}) rectangle ({(\j)*1.1-0.05},{(\i)*0.6+0.5});
      \node[font=\tiny] at ({(\j-0.5)*1.1-0.025},{(\i)*0.6+0.25}) {Tr};
    \fi
  }
  \node[font=\tiny,left] at (0,{(\i)*0.6+0.25}) {Fold \i};
}
\end{tikzpicture}
\end{center}
\vspace{0.2em}
\begin{exampleblock}{Vantagem}
  Usa \textbf{todos os dados} para treino e validação.
\end{exampleblock}
\end{columns}
\end{frame}

% ==============================================================================
% REFERÊNCIAS
% ==============================================================================
\begin{frame}{Referências}
\begin{thebibliography}{9}
\setbeamertemplate{bibliography item}[text]
{\small
\bibitem{quinlan1986} Quinlan, J. R. (1986). Induction of decision trees. \textit{Machine Learning}, 1(1), 81--106.
\bibitem{breiman2001} Breiman, L. (2001). Random forests. \textit{Machine Learning}, 45(1), 5--32.
\bibitem{vapnik1995} Vapnik, V. N. (1995). \textit{The Nature of Statistical Learning Theory}. Springer.
\bibitem{chen2016} Chen, T., \& Guestrin, C. (2016). XGBoost. \textit{KDD 2016}, 785--794.
\bibitem{ke2017} Ke, G., et al. (2017). LightGBM. \textit{NeurIPS 2017}, 3149--3157.
}
\end{thebibliography}
\end{frame}

\begin{frame}[c]
\begin{center}
{\Large \textbf{Consulte o Material Didático}}\\[0.8em]
{\large \textcolor{mineblue}{Módulo 3: Técnicas de Classificação}}\\[1.2em]
{\normalsize Exemplos práticos, exercícios resolvidos e implementações em Python disponíveis no \textbf{Notebook Colab}.}\\[1.5em]
\faGithub\ \href{https://github.com/eltonsarmanho}{github.com/eltonsarmanho}\\[0.4em]
\faEnvelope\ \href{mailto:eltonss@ufpa.br}{eltonss@ufpa.br}
\end{center}
\end{frame}

\end{document}
